
% =====================================================================
\section{Evidence ledger, search and update record}
\label{app:ledger}
The ledger, provided as ancillary files, consists of eight linked tables with stable identifiers.
\begin{itemize}
\item \texttt{entries.csv} (\nEntries{} rows): comparison level, model, observable, kind of match, timeline eligibility, dates with precision flags, elapsed days; the fields of the comparison record (instance and version, boundaries, tolerance, meaning of the failure probability, workload, access model); and three status fields, one of which records the basis on which a classical result is validated (a measured run, a benchmark evaluation, an experimental confirmation, an empirical extrapolation, a scaling law or a theorem).
\item \texttt{pairs.csv} (\nPairs{} rows): the same comparison record for the headline comparisons that are problem classes and not entries.
\item \texttt{resources.csv} (\nResources{} rows): one row per Hamiltonian, algorithm, layout and schedule, with source version, locator and locator level, the artifact read and the day it was retrieved, gate type, count and its scope, width and width convention, physical qubits, expected and one-shot runtime with the logical failure probability, space--time volume, coverage, unvalidated assumptions and our own extraction gaps.
\item \texttt{values.csv} (\nValues{} rows): classical times and other scalar facts with unit, kind, machine or criterion, version, locator and artifact.
\item \texttt{models.csv}: parameters of the formula figures and of the repetition scenario.
\item \texttt{forward\_classes.csv} (\nClasses{} rows) and \texttt{figure\_points.csv} (\nPoints{} rows; for each plotted point, the record, field and named transformation that each coordinate is computed from).
\item \texttt{update\_log.csv} (\nLog{} rows): each source check, with the version inspected, any access issue, and a marker on checks that a later one superseded.
\end{itemize}
Missing values carry one of four codes: \texttt{not\_reported} (the source does not give it), \texttt{not\_extracted} (we have not read it out), \texttt{not\_applicable}, \texttt{not\_estimated}.

\paragraph{Build and validation.} \texttt{build\_ledger.py} is the authoritative source and writes the tables. \texttt{make\_figdata.py} then runs \nChecks{} checks and writes the plot data only if all pass.
\begin{itemize}
\item Each of the \nCoords{} plotted coordinates is recomputed from its cited record and transformation and compared with the stored coordinate.
\item A point is refused if its record is known only from an abstract or at second hand.
\item A point is refused unless its source version is an arXiv identifier with a version number or a journal reference. Strings such as ``latest'' or ``as read for an earlier draft'' are rejected.
\item Units are checked: each transformation declares the unit it accepts and returns, and the result must be the unit of the axis. T-state and Toffoli counts cannot share a series.
\item A point's bound kind must equal that of its source row, and a one-sided value may be written only to a series drawn with a limit marker. No plotted value in this version is one-sided.
\item Every row taken from the dense-code table must satisfy that table's two internal relations: volume equals qubits times expected runtime, and expected runtime equals one-shot runtime divided by one minus the failure probability.
\item Links, vocabularies and the coverage-aggregation rule are checked.
\end{itemize}
\texttt{test\_negative.py} corrupts a copy of the ledger in fourteen ways, one for each kind of check, and passes only if the validator rejects each. The validation establishes that figures agree with tables and that tables are internally consistent. It does not establish that a value was read correctly from its source, which rests on the locators, and it does not make the identity of a version certain: we record the version stamp of the artifact we read, because a request for an older version of a page can silently return the current one.

\paragraph{What is and is not generated.} All data points, connectors, range bars and ratios quoted in the text are generated from the tables. Figs.~\ref{fig:quadratic} and~\ref{fig:readout} are formula figures whose parameters are in \texttt{models.csv}. Fig.~\ref{fig:taxonomy} is a diagram. Axis labels, legend text, the names printed beside some points, the vendor target lines of Fig.~\ref{fig:resource} and guide lines are typed in the manuscript. The summary tables in the text and Table~\ref{tab:master} are typed and were checked against the ledger by hand.

\paragraph{Completeness of the comparison records.} Table~\ref{tab:complete} counts, for the \nEoneEtwo{} device-comparison and resource-estimate entries, how many have each field of the comparison record filled. Most fields are still empty for most entries, so the comparison record is a schema we apply and report against, not yet a complete description of each pair. Table~\ref{tab:pairs} gives the records for the headline comparisons, with what is missing from each.

\begin{table}[htbp]
\centering\small
\caption{Completeness of the comparison record for the \nEoneEtwo{} entries of types E1 and E2, generated from \texttt{entries.csv}.}\label{tab:complete}
\begin{tabular}{@{}lrr@{}}
\toprule
Field & Filled & Not extracted\\\midrule
\completenessrows
\bottomrule
\end{tabular}
\end{table}

{\scriptsize
\setlength{\tabcolsep}{3pt}
\begin{longtable}{@{}L{2.1cm} L{2.5cm} L{3.2cm} L{2.8cm} L{2.1cm} L{2.9cm}@{}}
\caption{Comparison records for the headline pairs, generated from \texttt{pairs.csv}.}\label{tab:pairs}\\
\toprule
Pair & Instance & Tolerance and its meaning & Failure probability & Workload & What is missing\\\midrule\endfirsthead
\toprule
Pair & Instance & Tolerance and its meaning & Failure probability & Workload & What is missing\\\midrule\endhead
\bottomrule\endfoot
Fermi-Hubbard 10x10 (and 6x6) & half filling, U/t=4, nearest-neighbour hopping & 0.01 in total energy (source units); classical time is for a variational state at that accuracy & not extracted & one energy & sign-free QMC cost; parallel or GPU DMRG time; preparation cost\\
J1-J2 Heisenberg 10x10 & J2/J1=0.5 & 0.01 in total energy; classical time is for a variational state & not extracted & one energy & classical time on a current machine; neural-state result on the same cylinder\\
FeMoco, 76-orbital model & Li et al. active space, 76 orbitals, 113 electrons & 1 kcal/mol; classical: empirical extrapolation range (0.31 consensus, 1.25 conservative); quantum: phase-estimation precision & quantum: logical failure 0.070 per shot at minimum volume, spectral-hit probability not reported; classical: none & one energy per candidate state & measured overlap; candidate search on both sides; classical run at D near 393,000\\
Cytochrome P450 compound I, model X & 58 orbitals, 63 electrons & chemical accuracy on one energy & quantum: logical failure 0.098 per shot at minimum volume; classical: none & at least two spin states for a gap & both spin states at a stated tolerance on either side; a classical gap converged in bond dimension\\
X-ray absorption, Li4Mn2O cluster & active space of 22 electrons in 18 orbitals & Lorentzian broadening eta=0.05 Ha; spectral error by the source's perturbative analysis & not extracted & 2500 shots over 200 time points & validated initial state; classical spectrum for this cluster; a machine with the assumed clock\\
OTOC(2), 65 qubits & Willow processor, 65 of 105 qubits & signal-to-noise ratio of the measured correlator against exact simulation at up to 40 qubits (3.9 for the off-diagonal correlator); at 65 qubits accuracy is estimated from an empirical error model & none stated as a probability: error bars come from an empirical error model & quantum: all experiments (about 2 h); classical: one data point (3.2 y est.) & a common workload basis; tolerance; a converged classical method\\
Doped Clifford sampling (two instances) & v1: 70 q, 70 layers, 468 T; v2/v3: 64 q, depth 73, 314 T & v1: log-XEB 0.350 [0.298, 0.403]; v3: fidelity lower bound 0.349 & 95\% interval or confidence & v1: 2051 amplitude batches & a classical response to the periodic instance; fresh samples from either distribution\\

\end{longtable}
}

\subsection*{Search, inclusion and update record}
\label{app:search}
\textbf{Sources.} arXiv (quant-ph, cond-mat, physics.chem-ph, physics.flu-dyn, hep-lat, cs.LG), Nature and Science family journals, APS and ACS journals, vendor technical releases and the Quantum Advantage Tracker, 2017 to 11 November 2026.

\textbf{Inclusion.} A quantum hardware result or fault-tolerant resource estimate and a classical or AI result addressing the same or a closely related task, or a theoretical result bearing directly on such a pair. Context entries were added where an AI result is commonly cited as making quantum computing unnecessary. Timeline eligibility is defined in Fig.~\ref{fig:timeline}. Not entered in the timeline: later random-circuit and boson-sampling generations individually, the protein-fragment study (it claims model accuracy, not intractability), and the hybrid chemistry workflows (no advantage claimed).

\textbf{Exclusion.} Quantum results without any classical comparison; classical results without a quantum counterpart or stated quantum motivation; news reports without a primary source. For the standing claims we found no full classical reproduction; that is a search result as of the cutoff, not proof that none exists.

\textbf{Limitations.} The original searches were not logged query by query, so the selection cannot be reproduced exactly. Source checks made in November 2026 are logged in \texttt{update\_log.csv} (Table~\ref{tab:update}).

{\scriptsize
\setlength{\tabcolsep}{3pt}
\begin{longtable}{@{}L{1.25cm} L{2.9cm} L{2.6cm} L{5.0cm} L{4.0cm}@{}}
\caption{Source checks in November 2026, generated from \texttt{update\_log.csv}. Dates are JST. A check that a later one replaced is kept and marked SUPERSEDED, with the reason.}\label{tab:update}\\
\toprule
Date & Source & Check & Result & Version; issue\\\midrule\endfirsthead
\toprule
Date & Source & Check & Result & Version; issue\\\midrule\endhead
\bottomrule\endfoot
11-09 & microsoft.\allowbreak github.\allowbreak io/\allowbreak skala & Skala 1.0/1.1 WTMAD-2 & 3.89 / 2.72 & model cards\\
11-09 & arXiv:2506.\allowbreak 15784 & XAS gate unit & v2: T gates; current version: Toffoli gates per circuit & v2 abstract; current PDF abstract; SUPERSEDED on 10 and 11 Nov: the v2 PDF states Toffoli; the abstract page of v2 stated T gates\\
11-09 & arXiv:2210.\allowbreak 14109 v3 & Hubbard setting & Half filling; eps total energy; QPE $<$1e5 s at 10x10 single factory; DMRG Table S2 & v3; SUPERSEDED on 10 Nov by the Table S12 schedules\\
11-09 & arXiv:2605.\allowbreak 30455 & FeMoco-54 frontier & 36x space, 6.6x space-time vs SA; table values; P450 Cpd1X target listed & v1; SUPERSEDED on 11 Nov by the Table XXVII rows read from the PDF\\
11-09 & arXiv:2511.\allowbreak 02125 & Helios pairing paper metadata & Granet et al., v3 & v3\\
11-09 & arXiv:2607.\allowbreak 25998 & OLE primary & Barron et al. & v1\\
11-09 & arXiv:2603.\allowbreak 29091 & Gauge estimates & Table of fiducial SU(3) estimates & v3; subgroup value 3.5e12 not re-read; removed from figures in v6\\
11-09 & PMC9499570 & P450 DMRG costs & Model X costs estimated; model G measured & published\\
11-09 & research.\allowbreak google blog & OTOC(2) timing & $\sim$2 h experiments; 3.2 y per data point; 13,000x & blog; no calibration/readout statement; workload bases differ\\
11-09 & pnnl.\allowbreak gov; arXiv:2401.\allowbreak 04070 & Electrolyte validation & JACS 146, 20009 (2024); synthesis and conductivity & published\\
11-09 & arXiv:2609.\allowbreak 12146 & FeMoco-76 SA count & 1459 logical qubits, 9.99e8 Toffoli per PE run (citing Low et al.) & v1; secondary tabulation; excluded from figures in v6\\
11-09 & arXiv:2409.\allowbreak 11748 & FeMoco candidate-refinement protocol & 7.3e10 Toffoli total & v1 abstract; abstract-level; excluded from figures and ratios\\
11-10 & arXiv:2210.\allowbreak 14109 & Lattice QPE schedules & Table S12: single-factory 1.23e4 s (Hubbard 6x6), 1.05e5 s (Hubbard 10x10), 5.64e4 s (Heisenberg 10x10); (16,16) schedules 1.59e3, 1.46e4, 4.32e3 s & v3 HTML; supersedes the rounded $<$1e5 s statement used in v5\\
11-10 & arXiv:2210.\allowbreak 14109 & Lattice DMRG baselines & Tables S1, S2: single AMD EPYC 7702, 4 cores; time to a state within 0.01 of extrapolated E0; cylinders; Heisenberg 10x10 J2=0.5: 7.9e4 s & v3 HTML; text and Fig. S1 caption disagree on which direction is periodic\\
11-10 & arXiv:2506.\allowbreak 15784 & XAS resources & Table I N=18: 3.11e8 Toffoli longest circuit, 4.58e11 algorithm; 100 logical qubits minimum; 350 active-volume; 8.1e10 cycles; D=1e4; S=2500; RAS $\sim$32 d for 1000 states & v2 (20 Aug 2025); SUPERSEDED on 11 Nov by v3. The note written here on 10 Nov, that the v2 abstract states Toffoli, was wrong for the abstract web page (T gates) and right for the PDF\\
11-10 & arXiv:2605.\allowbreak 30455 & Dense-code compilations & Fig. 1 table (FeMoco-54); Table XXVII rows for FeMoco-54, FeMoco-76, Cpd1X 58o63e, XVIII 56o64e & v1; SUPERSEDED on 11 Nov: the 2.692 M-qubit FeMoco-76 row was dropped because a merged cell failed the volume check\\
11-10 & arXiv:2401.\allowbreak 07184 & QAC first posting & v1 14 January 2024 & arXiv listing; v5 used 12 January, inferred from the identifier\\
11-10 & arXiv:2607.\allowbreak 24937 & Heron Floquet & up to 74 q; TN not converged; Pauli-path truncation-dependent; RIKEN co-authors & v1 abstract\\
11-10 & arXiv:2607.\allowbreak 25941 & Doped Clifford versions & v1 28 Jul 2026; v2 1 Sep 2026; v3 2 Sep 2026 (64 q, depth 73, 314 T, 76 physical) & arXiv listing; v3 abstract; Note added not read\\
11-10 & arXiv:2608.\allowbreak 13110 & Doped Clifford classical response & 37.3 min, 32 nodes x 8 H100, 2051 amplitude batches; 10.6 min is a projection & v1 (13 Aug 2026) abstract; later version not read\\
11-10 & arXiv:2210.\allowbreak 08556 & Ultracold-atom sampling & 64 sites, 20 atoms; sampling-rate claim; Bayesian tests & listing and abstract; SUPERSEDED on 11 Nov: the request for v1 returned the current abstract, and the claim was dated from v1\\
11-10 & arXiv:2202.\allowbreak 09372; arXiv:2307.\allowbreak 09442 & Rydberg MIS pair & 18 Feb 2022; 18 Jul 2023 & arXiv listings and abstracts\\
11-10 & arXiv:2605.\allowbreak 01138 & Embedded protein-ligand workflow & 239 QPU-h; 21,006 circuits; 4.45e6 Fugaku node-h; no advantage claimed; no sampler ablation & PDF; abstract text garbled in extraction\\
11-10 & arXiv:2603.\allowbreak 28648 & GPU DMRG Fe4S4 & E=-327.24466 Ha at D\_SU(2)=12288; 12.6 h on one DGX B200 node & v1 HTML; instance match to the SQD experiment not established\\
11-10 & arXiv:2310.\allowbreak 05715 & Transformer NQS J1-J2 & E/N=-0.497634(1); 20 A100 GPUs, 4 days; DMRG -0.495530 & v2 HTML; boundary conditions not stated in retrieved text\\
11-10 & arXiv:2507.\allowbreak 02644 & NQS doped Hubbard & abstract only & v2 listing; no numbers extracted\\
11-11 & arXiv:2506.\allowbreak 15784 & XAS resources, version 3 & Table 1 N=18: 3.54e8 Toffoli (active volume 1.89e10) longest circuit; 5.20e11 Toffoli (2.76e13) algorithm; 100 logical qubits minimum; 350 active-volume; 1.08e8 and 1.58e11 logical cycles; less than two days at MHz & v3 (5 Nov 2026), PDF text layer; classical RAS estimate not re-read in v3\\
11-11 & arXiv:2605.\allowbreak 30455 & Dense-code Table XXVII & Full blocks for FeMoco-54, FeMoco-76, Cpd1X 58o63e; selected rows for XVIII 56o64e; columns qubits, volume, expected and one-shot runtime, p\_fail, Toffoli; Sec. VI.B defines expected runtime & v1, PDF text layer; first row of the XVIII 56o64e block not re-read\\
11-11 & arXiv:2210.\allowbreak 08556 & Ultracold-atom versions & v1 16 Oct 2022: another title, 32-site chain, power comparable to the fastest supercomputer; v2 11 Apr 2026: 64 sites, 20 atoms, three orders of magnitude; PRX 16, 021057 & abstract pages of v1 and v2\\
11-11 & arXiv:2403.\allowbreak 00910; 2503.\allowbreak 05693; 2505.\allowbreak 22514 & First-posting dates & 1 Mar 2024; 7 Mar 2025; 28 May 2025 & arXiv listings; Tindall et al. now Science 392, 868 (2026)\\
11-11 & arXiv:2607.\allowbreak 24937 & Heron Floquet validation & Unbiased and extrapolated estimators agree over their common window; cross-platform checks at selected cycles on two trapped-ion systems; thermodynamic-limit statement is an extrapolation & v1, PDF text layer\\
11-11 & arXiv:2607.\allowbreak 25941 & Doped Clifford revision & v3 text: periodic boundary conditions eliminate the low-width contraction path used in the classical response; framework unchanged & v3, PDF text layer\\
11-11 & arXiv:2608.\allowbreak 13110 & Doped Clifford classical response & v2 posted 29 Oct 2026; abstract figures unchanged (37.3 min; 32 nodes x 8 H100; log-XEB 0.35034 [0.29763, 0.40305]) & abstract page of v2\\
11-11 & arXiv:2605.\allowbreak 01138 & Embedded protein-ligand workflow & v2 5 Sep 2026: 21,006 circuits, over 239 hours, up to 94 qubits & abstract page of v2\\
11-11 & arXiv:2211.\allowbreak 07629 & Beverland et al. tables & Table I: 230/2740/25481 logical qubits, 2.4e6/5.4e11/1.5e10 T; Table IV runtimes and physical qubits for four qubit scenarios & v1, PDF text layer; pins the rows carried as (as read for v4)\\
11-11 & arXiv:2011.\allowbreak 03494 & THC FeMoco tables & Table IV (Reiher, 54 orbitals, M=350): 5.3e9 Toffoli, 2142 logical qubits; Table V (Li, 76 orbitals, M=450): 3.2e10 Toffoli, 2196 logical qubits; layout (Reiher only): 4e6 physical qubits, 6.7e9 Toffoli in 3 days at 25 kHz & v3, PDF text layer; CORRECTION: versions 4 to 6 attributed the Reiher numbers to the 76-orbital model\\
11-11 & arXiv:2501.\allowbreak 07231 & QSCI limits & 4,260 vs 50,551 determinants at 0.1 Ha; 0.1 and 0.0005 determinants per sample at 1e3 and 1e9 samples; 4e14 samples for micro-hartree precision & v5, PDF text layer; pins the rows carried as (as read for v4/v5)\\
11-11 & arXiv:1505.\allowbreak 06552; 1512.\allowbreak 05903; 2004.\allowbreak 06516; 2011.\allowbreak 03185; 2307.\allowbreak 09593; 2211.\allowbreak 05246; 2406.\allowbreak 06323; 2509.\allowbreak 08807; 2502.\allowbreak 20908; 2606.\allowbreak 01067; 2303.\allowbreak 13012 & Linear solvers and CFD & Complexity limits and end-to-end estimates as cited in Section 4.5 & abstract pages and search results; SUPERSEDED later on 11 Nov (v7.1) for 1505.06552, 1512.05903, 2406.06323 and 2509.08807, which were read in the PDF\\
11-11 & arXiv:2406.\allowbreak 06323 & Drag estimate: instances and classical reference & Sphere at Re=1e1 to 1e8 and three ship hulls; classical reference is steady RANS on one workstation: 1.1 h (Re=1e3) to 10 h (Re=1e7) with meshing, Re=1e8 out of memory (Table 3); estimates are stated as lower bounds; inlet/outlet conditions not included & v3, PDF text layer (browser); v7.1; the product range 1e21 to 1e39 remains an abstract-level value\\
11-11 & arXiv:2509.\allowbreak 08807 & Navier-Stokes crossover & 2\^{}80 grid; 8.71e6 physical qubits at error 5e-4; 42.6 days; classical 130 years; stated speedup 1100x; output by state tomography of a spectrally compact state & v1, PDF text layer (browser); v7.1\\
11-11 & arXiv:1505.\allowbreak 06552 & EM scattering resource table & Table II, N=332,020,680: depth 1.8e29 with oracles and 3.30e25 without; T-depth 8.2e28 and 1.28e25; measurements 2.8e27 and 8.23e21 & v2, PDF text layer (browser); v7.1\\
11-11 & arXiv:1512.\allowbreak 05903 & FEM complexity table & Table I, optimal preconditioning: classical (|u|\_2/eps)\^{}(d/2), quantum ||u||\_1/eps, up to logarithms & v2, PDF text layer (browser); v7.1\\
11-11 & arXiv:2502.\allowbreak 15882 & Spectrum amplification counts & Table I: FeMoco-54 1137 qubits, 3.41e8 Toffoli; FeMoco-76 1459 qubits, 9.99e8 Toffoli; prior THC 2142/5.3e9 and 2196/3.2e10; symmetry-compressed THC 1512/4.3e9 for the 76-orbital model & v1, PDF text layer (browser); v7.1; replaces the second-hand record; confirms the v7 correction of the THC attribution\\
11-11 & arXiv:2409.\allowbreak 11748 & FeMoco candidate refinement & 7.3e10 Toffoli; 2 or 3 QPE iterations; bond dimension 4000; overlap squared about 0.9 (extrapolated); 2.3 times the single-shot THC estimate & v1, PDF text layer (browser); v7.1\\
11-11 & arXiv:2605.\allowbreak 30455 & Dense code: Ru minimum-qubit row; Fig. 1 table & XVIII-56o64e row 1: 75.91k qubits, 3200 h expected, 351 h per shot, p\_fail 0.890, 5.55e9 Toffoli (3.20 + 1.69 + 0.65 by part); Fig. 1 table agrees with the four FeMoco-54 rows of the ledger & v1, PDF text layer (browser); v7.1; closes the row marked not re-read\\
11-11 & arXiv:2506.\allowbreak 15784 & XAS classical estimate & Sec. 5: 0.76 h x 1000 states, about 32 days; called likely optimistic on the classical side; accepted in Quantum on 23 Sep 2026 & v3, PDF text layer (browser); v7.1; the estimate is now pinned to v3\\
11-11 & arXiv:2306.\allowbreak 03709 & Gaussian boson sampling: datasets covered & Borealis (M16, M72, M216, M288), Jiuzhang 2.0 and Jiuzhang 3.0 & v2, PDF text layer (browser); v7.1; the classical result covers these datasets, so entry 2 is a family record. CORRECTION: versions up to 7 dated this entry from the 2020 Jiuzhang experiment, which the paired classical result does not simulate; the claim date is now that of Jiuzhang 2.0 (arXiv:2106.15534v1, 29 Jun 2021)\\
11-11 & arXiv:2506.\allowbreak 10191 & OTOC timing basis & Fig. 4c: estimated time to compute the correlator of a single circuit on Frontier by tensor-network contraction & v1, PDF text layer (browser); v7.1; the classical basis is one circuit; the 2 h quantum figure is still known from the blog only\\
11-11 & arXiv:2506.\allowbreak 10191 & OTOC timing & Main text beside Fig. 4(c): 3.2 years on Frontier is about 13,000 times the data collection time of 2.1 hours per circuit & v1, PDF text layer (browser); v8 CORRECTION: version 7.1 said the quantum figure was known from the vendor post only and covered a set of experiments\\
11-11 & github.\allowbreak com/\allowbreak isi-usc-edu/\allowbreak qsub & Drag estimate: per-instance records & All eight structured sphere files opened at commit cc8d3ce: products 5.3e21 (Re=10) to 7.9e40 (Re=1e8); task string 0.05 failure tolerance; least-squares slope 2.745 against the published 2.68 & repository web view; v8; UNRESOLVED against the published range 1e21 to 1e39; source authors not contacted (author decision)\\
11-11 & arXiv:2509.\allowbreak 08807 & Navier-Stokes Appendix F & F12 181 logical qubits; F15 28,301 queries; F20 9.41e7 Toffoli and 3.94e8 rotations; F24 distance 25; F26 1,000 samples; F30 3.68e6 s; F31 226,250 and F33 8.16e6 physical qubits; F36-F38 classical operation counts and rate & v1, PDF text layer (browser); v8; 130 years is the banded-Cholesky model\\
11-11 & arXiv:2511.\allowbreak 17437; 2608.\allowbreak 17115 & Gauge counts & Osorio: 3.5e12 T (Kogut-Susskind), 7.1e12 T (improved), 9-qubit register. Draper Table 5: U\_p 1.2e5 (B=6), 1.2e6 (B=9) T; Table 6 register terms & v1 of each, PDF text layer (browser); v8 CORRECTION: data width 27,000, not 24,000\\
11-11 & arXiv:2512.\allowbreak 03758; 2502.\allowbreak 11239; 2502.\allowbreak 20908; 2606.\allowbreak 01067 & Solver studies & Jennings Eqs. 181-186 and Re\^{}0.287; Tu: error 1e-5, 10 ns logical cycle, no input/output cost; Lapworth: 879/8,348 against 249/41,181; Goldfriend Table 2 and factor 2 in outer iterations & PDF text layer (browser); v8\\
11-11 & arXiv:2109.\allowbreak 11525; 2207.\allowbreak 10058; 2304.\allowbreak 12240; Nature 606, 75 & Boson sampling by experiment & Villalonga v1 (23 Sep 2021) covers the 2020 100-mode data and Jiuzhang 2.0; Martinez-Cifuentes v1 20 Jul 2022 (abstract); Jiuzhang 3.0 posted 24 Apr 2023; Borealis published 1 Jun 2022 & PDF text layer (v1) and listing pages; v8 CORRECTION: entry 2 split into four records; Jiuzhang 2.0 first response is 86 days, not 707\\
11-11 & doi:10.\allowbreak 1103/\allowbreak pb2g-j9cw & Spectral amplification, article of record & Open-access page loads (title: ...by Spectral Amplification); FeMoco paragraph present; text extraction refused & journal page (browser); NOT READ: schedule row stays second-hand\\
11-10 & github.\allowbreak com/\allowbreak physical-superintelligence/\allowbreak aaronson-ambainis & Claimed proof of the Aaronson-Ambainis conjecture & statement file matches Conjecture 1.7; simulation bound (4T)\^{}8 eps\^{}-14 delta\^{}-8 in paper Sec. 1.3 & commit 2a6ceae; Lean proof not rebuilt by us; no peer review\\

\end{longtable}
}

